\section{Text2Mem：记忆操作的统一表达与状态验证}
\label{sec:text2mem}

Text2Mem\upcite{wang-etal-2026-text2mem}面向自然语言驱动的记忆控制，提出统一操作语言及其基准。用户要求更新、合并或调整记忆时，系统通过结构化规范表达目标范围与参数，并检查执行后的状态是否符合预期。

\subsection{研究问题与适用范围}

长期交互不仅需要检索历史信息，还需要控制记忆的优先级、有效期、合并和删除。例如，“暂时不要再提午餐”可能意味着降低显著性，也可能被误解为删除记录。如果动作、范围与时间未被明确表达，系统即使调用成功，也可能改变错误的记忆。论文因此在自然语言与具体后端之间引入结构化操作规范。

中间规范用记忆生命周期操作表达意图，再由后端实现其语义，为不同后端提供共同的操作接口。论文主实验采用SQL/SQLite参考环境。

\subsection{操作集合、Schema与执行流程}

\paragraph{十二个操作。}
论文将操作分为编码、存储和检索三阶段。编码阶段仅有Encode；存储阶段包括Update、Label、Promote、Demote、Merge、Split、Delete、Lock、Expire九个操作；检索阶段包含Retrieve、Summarize。Encode组织内容、语义、来源与治理属性，向量生成可以延迟。Promote与Demote调整显著性，Demote可以归档而不删除；Delete处理软/硬删除，Expire设置到期规则，三者具有不同语义。

\paragraph{五字段契约。}
操作由\texttt{stage}、\texttt{op}、\texttt{target}、\texttt{args}、\texttt{meta}描述，分别对应阶段、动作、目标、专属参数和执行上下文。不同操作的必填条件不同，例如Encode创建新记忆时不必先选择已有目标。对于需要定位已有条目的操作，\texttt{target}在\texttt{ids}、\texttt{filter}、\texttt{search}、\texttt{all}四种选择方式中择一；跨字段约束限制锁定、到期和大范围写入等行为。

\paragraph{Validator--Parser--Adapter。}
Validator检查结构完整性与跨字段约束；Parser把合法实例转换为类型化操作，归一化时间和优先级等参数；Adapter再映射为具体后端动作。执行返回\texttt{ExecutionResult}，记录受影响条目和状态变化。多步骤请求可表示为保留依赖顺序的\texttt{schema\_list}。该流程检查结构和治理约束；目标选择与意图理解则由执行后的状态检查评估。

\paragraph{教学案例：提高课设进度记忆的优先级。}
以请求“提高本周课设进度笔记的重要性，其他记录不变”为例，设目标条目为17，可用STO阶段的Promote，将\texttt{ids=[17]}的\texttt{weight}设为0.9。状态检查同时覆盖目标权重、原有内容和非目标条目，确认操作只改变了指定范围。

\subsection{基准、指标与主要结果}

基准以自然语言指令、前置操作建立的初态、参考Schema和SQL状态断言组成评测实例，覆盖中英文、直接与间接指令、单操作与工作流。模型包括GPT-4o、Qwen2.5-72B-Instruct、Claude-4-Sonnet、Gemini-2.5-Pro与GPT-5，使用相同自然语言指令与提示模板。

评测包含三个指标：SMA用编辑距离或余弦相似度衡量生成Schema与参考的相似程度；ESR统计通过验证且无运行错误地执行的实例比例；EMR统计满足的SQL状态断言占全部期望断言的比例：
\begin{equation}
    \mathrm{EMR}=\frac{\sum_i\sum_{a\in A_i}\mathbf{1}[a\text{满足}]}{\sum_i|A_i|}.
\end{equation}
其中$A_i$为第$i$个实例的断言集合。EMR以断言为统计单位，断言较多的实例对总体指标影响更大。

\begin{table}[!htbp]
    \centering
    \caption{Text2Mem总体评测指标}
    \label{tab:text2mem-main}
    \small
    \begin{tabular}{lrrrr}
        \toprule
        模型 & SMA$_{\rm lev}$ & SMA$_{\rm cos}$ & ESR & EMR \\
        \midrule
        GPT-4o & 0.6776 & 0.8690 & 0.9484 & 0.2532 \\
        Qwen2.5-72B-Instruct & 0.6195 & 0.8349 & 0.9510 & 0.2535 \\
        Claude-4-Sonnet & 0.6230 & 0.8475 & 0.9594 & 0.3132 \\
        Gemini-2.5-Pro & 0.6491 & 0.8560 & 0.9479 & 0.2897 \\
        GPT-5 & 0.6115 & 0.8396 & 0.9662 & 0.2826 \\
        \bottomrule
    \end{tabular}
\end{table}

表\ref{tab:text2mem-main}列出各模型的总体指标。ESR约为94.79\%--96.62\%，EMR约为25.32\%--31.32\%，显示操作普遍可以执行，但状态效果与预期仍有较大差距。最高ESR与最高EMR分别对应GPT-5和Claude-4-Sonnet。

复杂存储和工作流的表现较弱。五个模型的Split EMR均为0；Merge仅GPT-4o为0.0465，其余为0。GPT-5的单操作EMR为0.2930，工作流EMR为0.1190，而工作流ESR为1.0000。这表明多步骤操作中的状态正确性是该基准的主要难点。

现有实验主要比较模型、操作类型和任务分组。Validator、Parser与Adapter的独立作用，以及指令复杂度对结果的影响，还需要组件消融和受控实验分析。

\subsection{方法局限与课设启发}

统一规范将动作、范围和生命周期参数显式化，便于执行前检查与执行后追踪。将结构相似、执行状态和实际效果分别计量，有助于定位目标选择、操作规划和后端执行中的问题。

方法的局限包括复杂记忆语义缺少形式保证、真实后端覆盖有限，以及尚未评测端到端长期智能体表现。状态断言的覆盖也会影响EMR：只检查目标权重，可能遗漏非目标条目被误改的情况。实际应用需要同时考虑后端适配、长期效果和错误传播。

StructMem关注事件绑定、整合与问答检索，Text2Mem关注操作表达和状态核验。与Semantic XPath的目标定位、MetaMem的事实利用结合比较，可以将记忆系统的错误依次定位到访问、推理和执行环节。

课程设计可借鉴显式输入约束与效果检查相结合的做法。若开展记忆更新实验，可固定初态、请求和模型，分别记录格式合法性、执行成功率、断言满足率与整任务通过率，并检查非目标记录是否保持完整。

% 原文定位：原文第1--3页，图1；第5、7页，\S3.3、\S4.3；第4页表1，第11页附录A；第4--5页，\S3.2；第3页图2、第5--7页\S3.3--4.1及图3；第12页附录B.1.1；后两项为本文的检查建议；第6--8页，\S4；第7--8页，\S4.1、\S4.3；第8、15页；第9页Limitations。
